% from ega2-1-fr.tex
\oldpage[II]{144}
d'un préschéma $Y$, on appelle \emph{codimension} de $F$ dans $Y$ et on note
$\operatorname{codim}_YF$ la borne inférieure des entiers
$\dim(\mathcal{O}_z)$, où $z$ parcourt $F$.
\end{env}

\begin{corollary}[7.3.5]
\label{II.7.3.5-fr}
Soient $Y$ un préschéma localement noethérien réduit, $N$ l'ensemble des
points $y\in Y$ où $Y$ n'est pas régulier
(\hyperref[0.4.1.4-fr]{0, 4.1.4})~; on suppose que
$\operatorname{codim}_YN\geq2$. Soit $f:X\to Y$ un morphisme de type fini,
séparé et fermé et soit $g$ une $Y$-section rationnelle de $X$~; si $Y'$ est
l'ensemble des points de $Y$ où $g$ n'est pas définie (ensemble qui est
\emph{fermé} (\hyperref[I.7.2.1-fr]{I, 7.2.1})) on a
$\operatorname{codim}_YY'\geq2$.
\end{corollary}

Il suffit de prouver que $g$ est définie en tout point $z\in Y$ tel que
$\dim\mathcal{O}_z\leq1$. Si $\dim\mathcal{O}_z=0$, $z$ est point générique
d'une composante irréductible de $Y$ (\hyperref[I.1.1.14-fr]{I, 1.1.14}), donc
appartient à tout ouvert partout dense de $Y$, et en particulier au domaine de
définition de $g$. Supposons donc que $\dim\mathcal{O}_z=1$~;
$\mathcal{O}_z$ est alors par hypothèse un anneau local noethérien régulier,
donc (\hyperref[II.7.1.6-fr]{7.1.6}) un \emph{anneau de valuation discrète}.
Soit $Z=\operatorname{Spec}(\mathcal{O}_z)$~; comme $U=Y-Y'$ est partout
dense, $U\cap Z$ n'est pas vide (\hyperref[I.2.4.2-fr]{I, 2.4.2})~; soit $g'$
l'application rationnelle de $Z$ dans $X$ induite par $g$
(\hyperref[I.7.2.8-fr]{I, 7.2.8})~; il suffit de montrer que $g'$ est un
\emph{morphisme} (\hyperref[I.7.2.9-fr]{I, 7.2.9}). Or, $g'$ peut être
considérée comme une $Z$-section rationnelle du $Z$-préschéma
$f^{-1}(Z)=X\times_YZ$~; il est clair que le morphisme $f^{-1}(Z)\to Z$
correspondant à $f$ est fermé, et il résulte de
(\hyperref[I.5.5.1-fr]{I, 5.5.1, (i)}) qu'il est séparé~; on conclut alors de
(\hyperref[II.7.3.3-fr]{7.3.3}) que $g'$ est partout définie~; comme $Z$ est
réduit et $X$ séparé sur $Y$, $g'$ est un morphisme
(\hyperref[I.7.2.2-fr]{I, 7.2.2}).

\begin{corollary}[7.3.6]
\label{II.7.3.6-fr}
Soient $S$ un préschéma localement noethérien, $X$ et $Y$ deux
$S$-préschémas~; on suppose $Y$ réduit, et en outre que l'ensemble $N$ des
points $y\in Y$ où $Y$ n'est pas régulier soit tel que
$\operatorname{codim}_YN\geq2$~; on suppose enfin que le morphisme structural
$X\to S$ soit propre. Soit $f$ une $S$-application rationnelle de $Y$ dans
$X$ et soit $Y'$ l'ensemble des points de $Y$ où $f$ n'est pas définie~;
alors on a $\operatorname{codim}_YY'\geq2$.
\end{corollary}

On sait (\hyperref[I.7.1.2-fr]{I, 7.1.2}) qu'on peut identifier les
$S$-applications rationnelles de $Y$ dans $X$ aux $Y$-sections rationnelles
de $X\times_SY$~; comme le morphisme structural $X\times_SY\to Y$ est fermé
(\hyperref[II.5.4.1-fr]{5.4.1}), on peut appliquer
(\hyperref[II.7.3.5-fr]{7.3.5}), d'où le corollaire.

\begin{remark}[7.3.7]
\label{II.7.3.7-fr}
Les hypothèses faites sur $Y$ dans (\hyperref[II.7.3.5-fr]{7.3.5}) et
(\hyperref[II.7.3.6-fr]{7.3.6}) seront en particulier réalisées lorsque $Y$
est \emph{normal} (\hyperref[0.4.1.4-fr]{0, 4.1.4}), en vertu de
(\hyperref[II.7.1.6-fr]{7.1.6}).
\end{remark}

Nous pouvons caractériser les morphismes universellement fermés (resp.
propres) par une réciproque de (\hyperref[II.7.3.3-fr]{7.3.3})~:

\begin{theorem}[7.3.8]
\label{II.7.3.8-fr}
Soient $Y$ un préschéma (resp. un préschéma localement noethérien), $f:X\to Y$
un morphisme quasi-compact séparé (resp. de type fini). Les conditions
suivantes sont équivalentes~:
\begin{enumerate}
\item[\rm a)] $f$ est universellement fermé (resp. propre).
\item[\rm b)] Pour tout $Y$-schéma de la forme
$Y'=\operatorname{Spec}(A)$, où $A$ est un anneau de valuation (resp. un anneau
de valuation discrète) de corps des fractions $K$, l'application canonique
\[
\operatorname{Hom}_Y(Y',X)\longrightarrow
\operatorname{Hom}_Y(\operatorname{Spec}(K),X)
\]
correspondant à l'injection canonique $A\to K$, est surjective (resp.
bijective).
